\documentclass[10pt,a4paper]{amsart}
\usepackage[margin=19mm]{geometry}
\usepackage{amsmath,amssymb,mathtools}
\usepackage[scr=rsfs,cal=euler]{mathalfa}
\usepackage{xfrac}
\usepackage[unicode,hidelinks,bookmarks=false]{hyperref}
\newcommand{\ie}{i.e.\ }
\newcommand{\cf}{cf.\ }
\newcommand{\resp}{respectively\ }
\newcommand{\Subsection}{\subsection}
\providecommand{\nobreakdash}{\nobreak\hbox{-}}
\setlength{\emergencystretch}{2em}
\multlinegap0pt
\usepackage{bm}
\usepackage{bbm}
\usepackage{dsfont}
\usepackage{xspace}
\usepackage{cancel}
\usepackage{mathtools}
\usepackage{amsthm}
\makeatletter
\@ifundefined{Theorem}{\newtheorem{Theorem}{Theorem}}{}
\@ifundefined{Lemma}{\newtheorem{Lemma}[Theorem]{Lemma}}{}
\makeatother
\def\CF{{\mathcal{F}}}    \def\OS{{\mathcal{O}S}}
\def\CO{{\mathcal{O}}}
\def\IBr{\mathrm{IBr}}
\def\Irr{\mathrm{Irr}}           \def\tenK{\otimes_K}
\def\OQ{{\mathcal{O}Q}}
\def\Z{{\mathbb Z}}
\def\tenB{\otimes_B}
\def\tenO{\otimes_{\mathcal{O}}}
\def\ten{\otimes}
\title[On the p-part of the conductor of a generalised character]{On the p-part of the conductor of a generalised character}
\author{Markus Linckelmann}
\date{Source: arXiv:2604.01351v2 (CC BY 4.0)}
\begin{document}
\maketitle

\noindent\emph{Source and licence.} On the p-part of the conductor of a generalised character. Version arXiv:2604.01351v2 (CC BY 4.0), distributed under the Creative Commons Attribution 4.0 International licence, as recorded in the arXiv licence metadata for this exact version and shown on the abstract page https://arxiv.org/abs/2604.01351v2. Full rights notice: licenses/arxiv-2604.01351-CC-BY-4.0.txt.\par\smallskip

\noindent\textit{Editorial scope.} Counted unit: the Lemma whose source label is \texttt{lem8} in the article (source0.tex lines 899--906, source label \texttt{lem8}), delivered together with the article's own complete original proof of it (source0.tex lines 908--955, one \texttt{proof} environment of 48 lines). No mathematics is written, repaired or re-derived: statement and proof are the article's own text line for line. Required context retained uncounted: lem1 (lines 593--598); lem75 (lines 845--852). Every definition, display and label of those lines is retained unchanged. Omitted from this package and not redistributed: 1--592, 599--844, 853--898, 907--907, 956--1213. Nothing inside a retained excerpt is omitted or altered: every retained body is delivered exactly as the article's lines are written, so no omission receipt of any kind accompanies this package. Citations: the retained excerpts carry no citation command at all, so no bibliography entry of the article is transcribed and no reference is redistributed. Reference adaptation: none was needed and none was made; every reference inside the delivered unit resolves inside it, so no unresolved reference or citation is produced. Printed slips kept literal: no printed word, symbol, hypothesis or number of the counted statement or of its proof was altered. Layout and class: the article's own class is \texttt{amsart} with packages including amssymb, amsthm, amsmath, xy; the wrapper is the lane's pinned \texttt{amsart} editorial wrapper at 10pt on A4, which restates the theorem-like environments the retained text needs and adds the article's own macro definitions that the retained text needs (\texttt{\textbackslash CF}, \texttt{\textbackslash CO}, \texttt{\textbackslash IBr}, \texttt{\textbackslash Irr}, \texttt{\textbackslash OQ}, \texttt{\textbackslash Z}, \texttt{\textbackslash ten}, \texttt{\textbackslash tenB}, \texttt{\textbackslash tenO}), with extra wrapper packages (bm, bbm, dsfont, xspace, cancel, mathtools). The delivered numbering is the unit's own. Assets: no retained span contains a figure environment, a graphics-inclusion command, a PDF-page inclusion, a table or a TikZ picture; no image file is copied into this package, the package declares no asset dependency and no image file is redistributed with it. Licence: version arXiv:2604.01351v2 of this article is distributed under the Creative Commons Attribution 4.0 International licence, as recorded in the arXiv licence metadata for this exact version and shown on the abstract page https://arxiv.org/abs/2604.01351v2. A full rights notice is delivered at licenses/arxiv-2604.01351-CC-BY-4.0.txt. Characters: every retained excerpt is pure ASCII, so the delivered engine, XeLaTeX with its default Latin Modern text font, reports no missing character.
\begin{Lemma} \label{lem1}
With the notation above, let $\chi$, $\chi'\in$ $\Z\Irr(B)$ such that $\chi-\chi'\in$
$\Pr(B)$. Then, for any nontrivial $p$-element $u\in G$ and any $\varphi\in$
$\IBr(kC_G(\langle u\rangle)$ we have 
$d^u_{\chi, \varphi} = d^u_{\chi',\varphi}.$
\end{Lemma}

\begin{Lemma} \label{lem75}
Let  $Q$ be a fully $\CF$-centralised nontrivial subgroup of $P$.
For any primitive local idempotent $j'$ in $(A')^Q$ there is a primitive local
idempotent $j$ in $A^Q$ such that  the up to isomorphism 
unique indecomposable direct summand of the $\OQ$-$B$-module 
$V_Q\tenO jB$ with vertex $\Delta Q$ is isomorphic to the non-projective 
indecomposable direct summands of $j'M$.
\end{Lemma}

\begin{Lemma} \label{lem8}
With the notation above, assume that the character $\chi_V$ of $V$ as values
in $\Z$. Let $U$ be a finitely generated $\CO$-free
$B$-module.  Denote by $\chi$ the character of $U$ and by $\psi$ the
character of $M\tenB U$. Then, up to signs and up a suitable order,
the non-ordinary generalised deomposition numbers of $\chi$ and
$\psi$ coincide.
\end{Lemma}

\begin{proof}
We note first that
it suffices to show that every non-ordinary generalised decomposition
number of $\psi$ is up to sign  a non-ordinary generalised  
decomposition
number of $\chi$. Indeed, if this statement holds in general, then 
that same statement applied to the $B'$-module $M\tenB U$ and 
the $B$-$B'$-bimodule $M^*$ shows 
that every non-ordinary generalised decomposition number of $M^*\ten_{B'} M\tenB U$ is
also, up to a sign, a non-ordinary generalised decmposition number of
$M\tenB U$. Since $M^*\ten_{B'} M\tenB U\cong$ $U\oplus T$ for some projective
$B$-module $T$ it follows from Lemma \ref{lem1} that  the non-ordinary
generalised decomposition numbers of $U$ and of $M^*\ten_{B'} M\tenB U$
coincide, and hence up to indexing, coincide with those of $M\tenB U$ as well.

Let $u$ be a nontrivial $p$-element in $P$. Set $Q=\langle u\rangle$.
Since $\chi$ and $\psi$ are class functions, we may assume that $Q$ is
fully $\CF$-centralised.

Note that the character of any finitely generated $\CO$-free indecomposable 
$\OQ$-module with vertex a proper subgroup of $Q$ vanishes at $u$.
if $j$ be a primitive local idempotent in $A^Q$, then the generalised 
decomposition number of $\chi$ at $u$ and the irreducible Brauer character of $kC_G(u)$ 
determined by $j$ is equal to $\chi(uj)$. Similary, if $j'$ is a primitive local
idempotent in $(A')^Q$ then the generalised decomposition number  of
$\psi$ at $u$ and the irreducible  Brauer character  of $kC_{G'}(u)$ determined
by $j'$ is equal to $\psi(uj')$. 

Fix a primitive local idempotent $j'$ in $(A')^Q$. By the above, we need to 
show that there is a primitive local idempotent $j$ in $A^Q$ such that
$\chi(uj)=$ $\pm\psi(uj')$.

By Lemma \ref{lem75}  there is a primitive local
idempotent $j$ in $A^Q$ such that  the up to isomorphism 
unique indecomposable direct summand of the $\OQ$-$B$-module 
$V_Q\tenO jB$ with vertex $\Delta Q$ is isomorphic to the non-projective 
indecomposable direct summands of $j'M$.

Since the character of an $\OQ$-module vanishes at $u$ unless possibly
if $u$ belongs to a vertex of a summand of this module, it follows that
the characters of the left $\OQ$-modules $j'M\tenB U$ and
$V_Q\tenO jU$ coincide when evaluated at $u$. Moreover, the character of $V_Q$
evaluated at $u$ is in $\{1,-1\}$. (We use here the hypothesis that $V$  has a 
$\Z$-valued character.) Thus the evaluations at $u$ of the characters
of $V_Q\ten jU$ and $jU$ differ at most by a sign.
Together it follows that the generalised decomposition numbers
$\chi(uj)$ and $\psi(uj')$   differ at most by a sign. The result follows.
\end{proof}

\end{document}
